\documentclass[11pt]{article}
\usepackage[margin=25mm]{geometry}
\usepackage{fontspec,unicode-math}
\setmainfont{Latin Modern Roman}
\setmathfont{Latin Modern Math}
\usepackage{amsmath,amssymb,mathtools}
\usepackage{xurl,hyperref}
\hypersetup{colorlinks=true,urlcolor=blue,linkcolor=blue}
\setcounter{secnumdepth}{0}
\setlength{\parindent}{0pt}
\setlength{\parskip}{5pt}
\setlength{\emergencystretch}{3em}
\newcommand{\sref}[2]{\hyperlink{source:#1:#2}{[S#2]}}
\newcommand{\cataloguescope}[1]{\paragraph{Source question.}#1}
\newcommand{\floor}[1]{\left\lfloor#1\right\rfloor}
\begin{document}
\section{Individual Manipulation Power in Neutral Social Choice}\label{statement}

\subsection{Definitions and mathematical statement}

\paragraph{Shared notation.}Standard number systems are denoted by $\mathbb N,\mathbb Z,\mathbb Q,\mathbb R,\mathbb C$; index ranges, computational models and specialized notation follow the cited problem source and the definitions below.


Neutrality means that relabeling alternatives relabels the chosen outcome in the same way. A dictatorship always follows a designated voter\textquotesingle s top choice. Distance from dictatorship is the probability of disagreement under the prescribed profile distribution; $M_i(f)$ measures profitable manipulation by voter $i$.
\paragraph{Statement recorded in the supplied source.}
Let $f$ be a neutral social choice function for $n$ voters and $m$ alternatives, and let $M_i(f)$ be voter $i$'s manipulation power. For $m=3$, if $f$ is $\varepsilon$-far from dictatorship then $\sum_iM_i(f)\geq C\varepsilon^2$. Can the current $\Omega(1/n)$ lower bound on $\max_iM_i(f)$ be improved; can neutrality be replaced by $f$ being far from having range of size at most 2; and does an analogous theorem hold for $m>3$, perhaps with a bound polynomially weaker in $m$? \sref{statement}{2}
\paragraph{Scope qualification.}
The exact normalization of manipulation power and profile measure require the source. All three extension questions are retained.
\paragraph{Source status.}
The supplied catalogue labels this record open. Its source-status metadata was checked against the supplied snapshot on 2026-10-11. This import does not independently establish current open status or mathematical correctness. The catalogue literature assessment is dated 2026-08-21.
\subsection{Short English statement}
Improve individual-voter manipulation bounds, weaken neutrality, and extend the theorem beyond three alternatives.
\subsection{Sources}
\begin{itemize}
\item[\hypertarget{source:statement:1}{S1}] Ulam AI and the UnsolvedMath Contributors, UnsolvedMath: A Curated Collection of Open Mathematics Problems, catalogue entry. Catalogue curation and metadata: CC BY 4.0. The supplied source question is retained; definitions and the English summary are editorial additions.
\url{https://huggingface.co/datasets/ulamai/UnsolvedMath}.
\item[\hypertarget{source:statement:2}{S2}] Combinatorics (2008), Oberwolfach Reports, DOI 10.4171/owr/2008/01, pp. 5--78
\url{https://doi.org/10.4171/owr/2008/01}.
\item[\hypertarget{source:statement:3}{S3}] Mossel--Racz, quantitative theorem
\url{https://arxiv.org/abs/1110.5888}.
\item[\hypertarget{source:statement:4}{S4}] three-alternative quantitative theorem
\url{https://doi.org/10.1137/090756740}.
\textit{Underlying Oberwolfach Reports material retains its source-specific licence, including CC BY-SA 4.0 where applicable; see THIRD\_PARTY\_NOTICES.md.}
\end{itemize}
\paragraph{Preparation and review.}
Original question attributed to the cited source and UnsolvedMath contributors. Codex (GPT-6; exact serving revision not exposed) assisted with definitions, English summaries, source inspection and formatting. No new human mathematical authorship or independent proof verification is asserted. Catalogue identity checked against the complete supplied registry and existing repository statements; known flagged repetitions and obvious identical conjectures excluded. Full mathematical-equivalence/literature review has not been performed. Statement preparation and final review followed contribution prompts C and D; automated checks verify structure and identity, not open status.
\label{end:statement}
\end{document}
